\documentclass{article}
\usepackage{graphicx} % Required for inserting images
\usepackage{amsmath}
\title{Lesson 3 Homework - Calculus}
\author{William Wang}
\date{August 2026}

\begin{document}

\maketitle

\section{}
Let $\varepsilon>0$ be arbitrary. Then $\displaystyle\lim_{x\to 1}(4x-6)=-2$ means
there exists a $\delta>0$ such that
\[
0<|x-1|<\delta \;\implies\; |(4x-6)-(-2)|<\varepsilon.
\]
\section{}
We have
\[
|(4x-6)-(-2)| = |4x-4| = |4(x-1)| = 4|x-1|
\]

Setting $4|x-1| < \varepsilon$ gives
\[
|x-1| < \frac{\varepsilon}{4}
\]

\begin{align*}
\varepsilon = 0.1 \implies \delta = \frac{0.1}{4} = 0.025 \\
\varepsilon = 0.01 \implies \delta = \frac{0.01}{4} = 0.0025 \\
\varepsilon = 0.00001 \implies \delta = \frac{0.00001}{4} = 0.0000025
\end{align*}
\section{}
We have
\[
|(4x-6)-(-2)| < \varepsilon
\;\;\Longrightarrow\;\;
|4x-4| < \varepsilon
\;\;\Longrightarrow\;\;
4|x-1| < \varepsilon
\;\;\Longrightarrow\;\;
|x-1| < \frac{\varepsilon}{4}
\]

\[
\delta = \dfrac{\varepsilon}{4}
\]

\section{}
 Let $\varepsilon>0$ be arbitrary and set $\delta = \dfrac{\varepsilon}{4} > 0$.

So if we have that $0<|x-1|<\dfrac{\varepsilon}{4}$, we have
\[
|f(x)-L| = |(4x-6)-(-2)| = |4x-4| = |4(x-1)|
\]
\[
= 4|x-1| < 4\left(\frac{\varepsilon}{4}\right) = \varepsilon.
\]

So the definition of limit is satisfied and $\displaystyle\lim_{x\to 1}(4x-6) = -2$.

\end{document}
